% Generated draft: requires mathematical and source review.
\documentclass[11pt]{article}
\usepackage[margin=25mm]{geometry}
\usepackage{fontspec}
\setmainfont{Latin Modern Roman}
\usepackage{amsmath,amssymb,mathtools,unicode-math}
\setmathfont{Latin Modern Math}
\usepackage{enumitem,needspace,xurl,hyperref}
\setlength{\parindent}{0pt}
\setlength{\parskip}{5pt}
\setlength{\emergencystretch}{3em}
\newcommand{\N}{\mathbb{N}}
\newcommand{\Z}{\mathbb{Z}}
\newcommand{\Q}{\mathbb{Q}}
\newcommand{\R}{\mathbb{R}}
\newcommand{\C}{\mathbb{C}}
\newcommand{\F}{\mathbb{F}}
\newcommand{\A}{\mathbb{A}}
\newcommand{\PP}{\mathbb{P}}
\newcommand{\E}{\mathbb{E}}
\newcommand{\eps}{\varepsilon}
\newcommand{\dd}{\,\mathrm{d}}
\begin{document}
\section*{UnsolvedMath ID: 6700071}
\textbf{Source problem code:} catalogue entry
\subsection{Definitions and mathematical statement}
The proposed model is the Riemannian cylinder
\[
(S^{n-1}(1)\times\R,\ g_{\mathrm{rd}}+\mathrm dt^2),
\]
whose scalar curvature is \((n-1)(n-2)\). With \(m=n-1\), the conjectural implication in the source is
\[
\operatorname{macr.dim}(\Psi(X))\ge1,
\qquad \operatorname{Sc}_X\ge(n-1)(n-2)
\quad\Longrightarrow\quad X=S^{n-1}\times\R.
\]
The scalar inequality is pointwise. If equality of manifolds means metric rigidity, the conclusion should be an isometry to the displayed unit-cylinder metric; a diffeomorphism or homeomorphism conclusion would be weaker and different. The extraction does not explicitly select this meaning.

The construction \(\Psi(X)\) is missing, including its maps, domain/target, topology and metric. The research summary describes a degree-one map space, but supplies no definition of degree, admissible maps, or scale. Macroscopic dimension is a large-scale invariant whose precise controlled-map or bounded-cover convention must also be fixed; it is not the ordinary topological dimension of an unspecified function space.

The source begins if \(m=n-1\) and refers to this property, so earlier ambient hypotheses are essential. The supplied summary calls the manifold complete and nontrivial, but these terms, the dimension range, and any connectedness or additional geometric conditions are not included in the statement fragment. They are not invented from the cylinder example. The draft therefore records the exact conditional rigidity target, pending recovery of the space \(\Psi(X)\) and the initial class of \(X\), rather than claiming a theorem based solely on the scalar lower bound. The summary reports the conjecture unresolved.
\subsection{Short English statement}
Does positive macroscopic dimension of the intended map space, together with the sharp cylinder scalar-curvature bound, force the sphere-times-line cylinder? The map space, large-scale metric, ambient hypotheses, and meaning of equality are absent.
\subsection{Sources}
\begin{itemize}
\item Local UnsolvedMath registry, permanent ID \(6700071\), exact source code \texttt{catalogue entry}; printed scalar-curvature tag [?75], occurrence 74. Proposer Misha Gromov, 2017.
\item Gromov, \emph{101 Questions, Problems and Conjectures around Scalar Curvature} (2017), recorded PDF page 73; \url{https://www.ihes.fr/~gromov/wp-content/uploads/2018/08/101-problemsOct1-2017.pdf}. Recorded access 2026-07-29; pdf-text-normalized extraction.
\item Supplied OPEN-TRIAGE research summary describes a complete-manifold degree-one-map-space rigidity problem and reports no resolution, without additional precise definitions or publication references.
\end{itemize}
\end{document}
